\section{总结与个人反思}

本章的目的在于总结本次课程报告的主要工作，并反思在大数据分析流程中需要特别注意的问题。本文完成了从数据集选择、合规性检查、预处理、EDA、分类、回归、聚类到 LaTeX 报告组织的完整流程。

\subsection{主要结论}

第一，2025 Digital Lifestyle Benchmark Dataset 满足本课程报告对公开表格数据、样本规模、字段数量和多任务建模的要求。数据共有 3500 条记录和 24 个字段，但其合成数据属性决定了本文结论只能用于课程建模和行为分析练习。

第二，分类任务在严格排除目标泄漏字段后，最终采用 Gradient Boosting 和 threshold=0.14 的筛查策略，测试集 Recall=0.6420、F1=0.5355、PR-AUC=0.5084。该结果说明模型可以辅助识别一部分高风险样本，但不能作为诊断工具。

第三，回归任务显示 \texttt{digital\_dependence\_score} 更适合作为主线目标，最佳 $R^2=0.9839$、MSE=3.1471、MAE=0.9982；相比之下，\texttt{productivity\_score} 的最佳 $R^2=-0.0041$，属于弱预测/负结果。这个结果提醒我，实验报告不应只挑选效果好的目标，而应诚实呈现模型不能解决的问题。

第四，聚类任务在只使用数字行为与生活习惯数值特征的前提下得到三类探索性画像，但 KMeans $k=3$ 的 Silhouette=0.1860，说明分群边界不清晰。因此，聚类结果适合用于描述用户画像，不适合用于严格分类。

\subsection{课程收获}

通过本次报告，我对大数据分析流程有了更完整的理解。数据分析不只是运行模型，还包括数据来源审查、字段含义理解、缺失与重复检查、异常范围检查、特征工程、目标泄漏控制、评价指标选择和结果解释。尤其是在分类任务中，如果只看 Accuracy，就会忽视高风险样本召回不足的问题；在回归任务中，如果只报告效果好的数字依赖目标，就会掩盖生产力预测失败这一重要事实。

本次实验也让我认识到负结果的价值。\texttt{productivity\_score} 的预测效果较弱并不是报告失败，而是说明当前特征对该目标解释力有限。真实数据分析中，承认模型边界比给出过度结论更重要。

\subsection{改进方向}

后续可以从三个方向改进。第一，若能获得真实、匿名且合规的数字行为数据，可以检验当前模型在真实样本上的稳定性。第二，可以增加时间维度，例如连续多周的设备使用和睡眠记录，以分析行为变化而非单日静态状态。第三，聚类部分可以结合更多外部验证指标或专家解释，但在当前数据和 Silhouette 结果下，不宜继续强化分群结论。
